\documentclass[11pt]{article}

\usepackage[margin=1in]{geometry}
\usepackage{amsmath, amssymb}
\usepackage{listings}
\usepackage{xcolor}
\usepackage{hyperref}
\usepackage{enumitem}
\usepackage{titlesec}
\usepackage{parskip}
\usepackage{booktabs}

% Code listing style
\lstset{
    basicstyle=\ttfamily\small,
    keywordstyle=\color{blue}\bfseries,
    commentstyle=\color{gray}\itshape,
    stringstyle=\color{red},
    frame=single,
    breaklines=true,
    showstringspaces=false,
    tabsize=4
}

\title{\textbf{CS3210: Operating Systems} \\ Lecture 4 --- OS Abstractions: Isolation}
\author{John Kim \\ Georgia Institute of Technology}
\date{January 27, 2026}

\begin{document}

\maketitle
\tableofcontents
\newpage

% ============================================================
\section{Administrative Notes}
% ============================================================

\textbf{Lab 0} closed on Friday, January 23 with 148 submissions. If you submitted,
expect to see your grade on Gradescope shortly.

\textbf{Lab 1} (Startup, CPU Stacks, Debugging, and Physical Memory) was assigned on
January 23 and is due \textbf{February 6, 2026}. Lab 1 must be completed
\emph{individually}---no partners for this assignment. Use the full two weeks; the
autograder caps daily feedback submissions, so starting early is essential.

Starting with \textbf{Lab 2}, teams of 1--2 are permitted and encouraged. If you want a
partner, use the ``Search for Teammates!'' post pinned on Piazza---state your background,
availability, specialisation, goals, and the times of week you can pair-program. Student
groups will be created under Canvas $\to$ People once teams are finalised. Partner changes
are only allowed \emph{between} Labs 2, 3, and 4; once a lab begins you are committed
to your current partner until that lab closes.

% ============================================================
\section{What Is Isolation?}
% ============================================================

The word \emph{isolation} in everyday English means to set or place something apart---to
detach or separate it so that it stands alone. In the context of operating systems the same
intuition applies: \textbf{isolation} is the property by which one software entity (a
process, a user, or the kernel itself) is protected from the faults, misbehaviour, or
malicious intent of other entities sharing the same hardware.

This lecture develops that concept from first principles: what isolation gives us, what it
costs, when it is appropriate, and crucially, how an OS actually implements it on real
hardware.

% ============================================================
\section{The Case For Isolation: Benefits}
% ============================================================

Isolation serves three distinct and equally important purposes in OS design.

\textbf{Failure isolation} is the most fundamental. If one process crashes, scribbles
over memory, or enters an infinite loop, the damage must be contained. It must not cascade
into other processes, and above all it must not destroy the OS itself. This is one of the
primary motivations for the entire architectural split between user mode and kernel mode.
Prof.\ Tumanov will return to failure isolation extensively throughout the semester because
it is the single most important driver of kernel structure and design decisions.

\textbf{Performance isolation} ensures that one process's misbehavior cannot starve others.
A process that consumes all available CPU time, allocates every page of physical memory, or
holds disk I/O bandwidth hostage must be constrained. This isolation dimension raises a
fundamental policy question: should the OS offer \emph{best-effort} fairness (trying to be
fair but with no hard guarantees) or \emph{guaranteed} fairness (each process receives at
least a specified minimum share)? The answer depends on the workload and is a core
scheduling design decision.

\textbf{Correctness isolation} guarantees that the behaviour of one program is not affected
by the presence or activity of other programs running concurrently. Even if process B is
doing nothing harmful---it is just running---process A's outputs should be identical to what
they would be if A were the only program on the machine. Without correctness isolation, bugs
in process B can corrupt shared data structures or interfere with process A's memory in
subtle ways that are extremely difficult to diagnose.

% ============================================================
\section{The Cost of Isolation: Why Full Isolation Is Useless}
% ============================================================

Isolation has a dark side. Taken to its logical extreme, \emph{perfect} isolation makes a
process completely useless. A perfectly isolated process cannot communicate with anything:
it cannot write to the screen, receive keyboard input, send network packets, or even return
an exit code. The only scenario in which full isolation is acceptable is when the computation
is entirely self-contained---essentially a calculator that takes no input and produces no
observable output. That is almost never what we want.

The deeper insight is that isolation is not a binary property. The useful design space lies
between the two extremes. The key takeaway from this discussion is:

\begin{quote}
We want \emph{some} isolation \emph{sometimes}. How much isolation is appropriate depends on
the workload, the trust relationships between entities, and the current phase of execution.
These requirements can change dynamically.
\end{quote}

% ============================================================
\section{When Is Isolation Needed?}
% ============================================================

Two conditions together indicate that isolation is appropriate:

\begin{enumerate}
    \item \textbf{Shared resources.} Isolation is relevant whenever multiple entities
          access the same underlying physical resource. Physical memory is the canonical
          example: every process on a multiprocess OS is competing for the same finite pool
          of DRAM pages. Without isolation, process A can read or overwrite any physical
          address, including the memory of process B or the kernel itself.

    \item \textbf{Non-cooperative entities.} Isolation is especially important when the
          entities sharing a resource are not designed to cooperate---they are written by
          different authors, have different goals, and may not trust each other. Processes
          are the classic example of non-cooperative entities. Threads within a single
          process, by contrast, \emph{share} the same address space intentionally and are
          presumed to cooperate; imposing strict memory isolation between threads would
          defeat the purpose of threading.
\end{enumerate}

A useful diagnostic framework is to ask: \emph{What are the entities? What is the shared
resource?} Applying this framework reveals whether isolation is needed and at what
granularity it should be enforced.

% ============================================================
\section{Isolation Tradeoffs}
% ============================================================

Increasing isolation strength moves safety and usability in opposite directions:

\begin{center}
\begin{tabular}{lll}
\toprule
\textbf{Property} & \textbf{Weaker Isolation} & \textbf{Stronger Isolation} \\
\midrule
Safety     & Worse  & Better \\
Usability  & Better & Worse  \\
Performance & Unpredictable & Unpredictable \\
\bottomrule
\end{tabular}
\end{center}

The safety/usability tradeoff traces a \emph{Pareto frontier}: any point on the frontier is
a legitimate design choice; the goal of good OS design is to push that frontier outward
toward higher safety \emph{and} higher usability simultaneously---that is the ``dream'' that
active research aims for. A bad design achieves neither safety nor usability---it sits
below the frontier through unnecessary overhead or unnecessary restrictions. The two ends
of the current frontier are \emph{Reality 1} (high safety, low usability, like a highly
locked-down microkernel) and \emph{Reality 2} (low safety, high usability, like early
MS-DOS).

The performance axis is marked with a question mark in the lecture because the relationship
is genuinely ambiguous: isolation mechanisms such as paging add overhead (TLB misses, mode
switches), but isolation can also \emph{improve} performance by preventing resource
contention and enabling better scheduling decisions.

% ============================================================
\section{What Might We Want to Isolate?}
% ============================================================

Armed with the necessity conditions and the tradeoff framework, we can enumerate concrete
isolation targets:

\begin{itemize}
    \item \textbf{Kernel space from user space} --- the kernel manages the hardware and
          enforces all the other isolation guarantees; if user code could freely overwrite
          kernel memory, every other isolation mechanism collapses. This is the highest-priority
          isolation boundary in any OS.
    \item \textbf{User processes from each other} --- no process should be able to read
          another process's memory, monopolise the CPU, exhaust shared file descriptors, or
          otherwise interfere with a co-resident process's operation.
    \item \textbf{Memory regions} --- even within the kernel, different regions (BIOS area,
          kernel text, kernel heap, per-CPU data) should have appropriate access controls so
          that a kernel bug does not inadvertently overwrite interrupt vectors or code pages.
    \item \textbf{Security-sensitive code from ordinary code} --- cryptographic key material,
          authentication logic, and other security-critical functions should be isolated from
          general application code that might be compromised or buggy.
\end{itemize}

% ============================================================
\section{Isolation Mechanisms}
% ============================================================

Modern operating systems on x86 hardware use five principal mechanisms to achieve
isolation. These mechanisms are layered and complementary; each addresses a different
dimension of the isolation problem.

\begin{enumerate}
    \item \textbf{Hardware protection rings} --- the CPU enforces privilege levels that
          restrict which instructions and which resources are accessible at each level.
    \item \textbf{Address spaces} --- virtual memory with page mapping gives every process
          its own isolated view of memory, preventing cross-process address access.
    \item \textbf{Time-slicing} --- the scheduler allocates CPU quanta to processes,
          ensuring no process can monopolise the processor.
    \item \textbf{System call interface} --- a controlled, vetted gateway through which user
          code can invoke kernel services without exposing kernel internals.
    \item \textbf{Permissions and privilege levels} --- UNIX-style user/group/other file
          permissions and UID/GID separation between users.
\end{enumerate}

The remainder of this lecture examines each mechanism in depth with concrete xv6 examples.

% ============================================================
\section{Hardware Isolation: Protection Rings}
% ============================================================

\subsection{The x86 Ring Model}

The x86 architecture defines four \textbf{protection rings} (privilege levels 0 through 3).
Ring~0 is the most privileged; Ring~3 is the least privileged. The original design allocated
different rings to different software layers:

\begin{center}
\begin{tabular}{lll}
\toprule
\textbf{Ring} & \textbf{Privilege Level} & \textbf{Intended Occupant} \\
\midrule
0 & Highest & Operating System Kernel \\
1 & High    & OS Services \\
2 & Medium  & OS Services \\
3 & Lowest  & User Applications \\
\bottomrule
\end{tabular}
\end{center}

In practice, virtually all modern operating systems---including xv6, Linux, and
Windows---use only rings 0 and 3. Rings 1 and 2 were intended for OS services that need
some privilege but not full kernel access; this design was never widely adopted because the
complexity of managing three distinct privilege levels outweighed the benefits. The resulting
binary ``kernel mode / user mode'' model is simple, robust, and sufficient.

\subsection{What Ring 0 Protects}

Running in ring~0 (also called \textbf{kernel mode}) grants access to CPU resources that
are explicitly protected from ring~3:

\begin{itemize}
    \item \textbf{Writes to \texttt{\%cr3}} --- CR3 holds the base physical address of the
          page directory. A process that could write \texttt{\%cr3} could replace its own
          page table with any other, instantly gaining access to any physical memory. Ring~0
          protection of CR3 is therefore the foundation of all address-space isolation.
    \item \textbf{Writes to \texttt{\%cs}} --- the low two bits of \texttt{\%cs} encode the
          \textbf{Current Privilege Level (CPL)}. If user code could write \texttt{\%cs}
          directly it could elevate itself to ring~0, bypassing all privilege checks.
    \item \textbf{Every memory read/write is checked against the privilege level} ---
          page table entries carry a \textbf{U/S} (User/Supervisor) bit. When CPL = 3 and
          the U/S bit of the target page is 0 (supervisor-only), the hardware raises a
          general protection fault. This is the primary mechanism preventing user code from
          reading or writing kernel memory.
    \item \textbf{I/O port accesses are privileged} --- the \texttt{IN} and \texttt{OUT}
          instructions that communicate with hardware devices (including disk controllers)
          are restricted to ring~0. An unprivileged process that could issue arbitrary
          \texttt{OUT} commands could overwrite disk sectors, corrupt filesystems, or
          reprogram the interrupt controller.
    \item \textbf{Control register accesses} --- writes to EFLAGS, CR4, and other control
          registers that govern CPU operating mode are ring~0-only.
\end{itemize}

\subsection{The Segment Selector and Current Privilege Level}

The CPU knows the currently executing privilege level through the \textbf{segment selector}
stored in \texttt{\%cs}. A segment selector is a 16-bit value with the following structure:

\begin{itemize}
    \item \textbf{Bits 15--3 (Index)} --- an index into the GDT or LDT.
    \item \textbf{Bit 2 (TI, Table Indicator)} --- 0 selects the GDT; 1 selects the LDT.
    \item \textbf{Bits 1--0 (RPL, Requested Privilege Level)} --- the privilege level
          being requested for this segment reference.
\end{itemize}

The \textbf{CPL} is defined as the two least-significant bits of \texttt{\%cs}. When
CPL = 0 the processor is in ring~0 (kernel mode); when CPL = 3 it is in ring~3 (user mode).
When the kernel switches to user mode it sets CPL = 3 by loading a user-mode code segment
selector; when it takes a trap or interrupt it restores CPL = 0 by loading a kernel-mode
selector. The CPU itself enforces that ring~3 code cannot raise CPL by writing \texttt{\%cs}
directly---only the hardware trap mechanism may change CPL.

\subsection{The Code-Segment Descriptor}

A segment descriptor is an 8-byte entry in the GDT that describes the segment's base,
limit, and permission flags. For a code segment the relevant flags are: \textbf{P} (Present,
must be 1 for a valid segment), \textbf{DPL} (Descriptor Privilege Level, 2 bits---the
minimum CPL required to use this segment), \textbf{G} (Granularity---0 means the limit is
in bytes; 1 means pages), \textbf{D/B} (Default operand size, 1 for 32-bit), \textbf{L}
(Long mode bit), \textbf{C} (Conforming), \textbf{R} (Readable), and \textbf{A} (Accessed).

The DPL field is the key isolation lever: a user-mode code segment has DPL = 3, while the
kernel's code segment has DPL = 0. As a concrete example, the value \texttt{0x009b} in
\texttt{\%cs} under QEMU decodes as:

\[
\texttt{0x009b} = \text{0000 0000 1001 1011}_2
\]

Bits 15--3 give a GDT index; bits 1--0 = $\texttt{11}_2$ = 3, confirming CPL = 3 (user mode).

\subsection{CPL in xv6: The GDT Setup}

xv6 maintains a per-CPU GDT configured in two stages. The \textbf{bootstrap GDT}
(in \texttt{bootasm.S}) is minimal---just a null entry, a flat code segment, and a flat data
segment---sufficient to get into protected mode. Once the main kernel initialises each CPU,
it replaces this with a more complete per-CPU GDT:

\begin{lstlisting}[language=C, caption={vm.c --- per-CPU GDT initialisation showing DPL distinction}]
c = &cpus[cpunum()];
c->gdt[SEG_KCODE] = SEG(STA_X|STA_R, 0, 0xffffffff, 0);        // DPL=0
c->gdt[SEG_KDATA] = SEG(STA_W,       0, 0xffffffff, 0);        // DPL=0
c->gdt[SEG_UCODE] = SEG(STA_X|STA_R, 0, 0xffffffff, DPL_USER); // DPL=3
c->gdt[SEG_UDATA] = SEG(STA_W,       0, 0xffffffff, DPL_USER); // DPL=3
// Per-CPU private data pointer in GS
c->gdt[SEG_KCPU]  = SEG(STA_W, &c->cpu, 8, 0);

lgdt(c->gdt, sizeof(c->gdt));
loadgs(SEG_KCPU << 3);
\end{lstlisting}

The four \texttt{SEG(...)} calls create flat segments (base = 0, limit = \texttt{0xffffffff})
that span the entire 4~GB address space. The isolation does \emph{not} come from segment
limits---all segments are identical in base and limit. It comes from the \textbf{DPL}:
kernel code and data segments carry DPL = 0, while user code and data segments carry DPL = 3
(\texttt{DPL\_USER}). The CPU prevents ring~3 code from loading a ring~0 segment selector
into \texttt{\%cs}, so user code is confined to the user code and data segments and cannot
access kernel-privilege instructions.

\subsection{Segment Selectors and Interrupt Handling}

The segment selector mechanism also governs who may invoke which interrupt handlers. Each
entry in the \textbf{Interrupt Descriptor Table (IDT)} specifies the target \texttt{CS} and
\texttt{EIP} for the handler, plus a DPL in the gate descriptor. When an interrupt fires,
the CPU checks whether CPL $\leq$ DPL of the gate descriptor:

\begin{itemize}
    \item If CPL $\leq$ DPL, the interrupt is permitted and the handler runs.
    \item If CPL $>$ DPL (user trying to fire a ring~0-only gate), the CPU raises a general
          protection fault.
\end{itemize}

This mechanism ensures that user code cannot directly invoke privileged interrupt
handlers---the only valid entry point into the kernel is through approved gates in the IDT.

% ============================================================
\section{The Unit of Isolation: The Process}
% ============================================================

The central unit of isolation in a traditional OS is the \textbf{process}. Four goals
follow from this design choice:

\begin{enumerate}
    \item \textbf{Process-to-process isolation} --- process~X must not be able to read
          or write process~Y's memory, exhaust Y's file descriptors, or consume CPU in a
          way that starves Y. This applies to memory, CPU time, open file descriptors, and
          any other limited resource.
    \item \textbf{Process-to-kernel isolation} --- a process must not be able to disable
          or circumvent the kernel's enforcement of isolation for other processes. If a
          process could crash the kernel, the protection guarantees for every other process
          collapse instantly.
    \item \textbf{Robustness against bugs and malice} --- isolation must hold even if a
          process contains bugs (which is inevitable) or is deliberately adversarial (which
          must be assumed). The infamous \textbf{Spectre} and \textbf{Meltdown} hardware
          vulnerabilities demonstrate that even well-designed software isolation can be
          undermined by hardware-level attacks that exploit speculative execution; these
          attacks are discussed in Section~\ref{sec:limits}.
    \item \textbf{Failure containment} --- if process~A has a bug, that bug must not
          propagate to process~B unless B is a child of A and the relationship was explicitly
          established. Independent processes must be truly independent.
\end{enumerate}

A subtle and interesting question: can we isolate a process \emph{from the kernel itself}?
Normally the kernel is trusted. But the research on microkernels and capability-based systems
explores exactly this question, attempting to build OSes where even kernel components are
isolated from each other.

% ============================================================
\section{Memory Isolation via Virtual Memory}
% ============================================================

\subsection{The Virtual Address Space Abstraction}

\textbf{Virtual memory} is the primary mechanism for memory isolation. The core idea is to
give every process the illusion that it has its own private address space spanning the
entire addressable range of the CPU---on 32-bit x86, that is 4~GB from address
\texttt{0x00000000} to \texttt{0xFFFFFFFF}---while the OS uses the hardware
\textbf{Memory Management Unit (MMU)} to map each virtual address to a distinct physical
address.

Because each process has its own page table (pointed to by CR3), a virtual address
\texttt{0x10000} in process~A and the same virtual address \texttt{0x10000} in process~B
map to \emph{completely different physical pages}. Process~A cannot read process~B's memory
by guessing virtual addresses, because A's page table simply does not contain mappings to
B's physical pages. This gives \emph{address space isolation} automatically: the mapping
hardware enforces it without any explicit per-access check by the kernel.

In xv6, the virtual address space is divided at \texttt{0x80000000}:

\begin{itemize}
    \item \textbf{Below \texttt{0x80000000}}: user space. Contains the user text and data,
          heap, and user stack.
    \item \textbf{Above \texttt{0x80000000}}: kernel space. The kernel's text and data
          (beginning at \texttt{0x80100000}) and free memory are mapped here, with the BIOS
          area placed at \texttt{0x80000000}.
\end{itemize}

The kernel mappings are \emph{present in every process's page table}---the same physical
kernel pages are mapped at the same high virtual addresses in all processes. This allows
the kernel to run using the same virtual addresses regardless of which process is currently
active. However, the kernel pages are marked \emph{supervisor-only} (U/S = 0), so user
code that tries to dereference a kernel address gets a page fault rather than a successful
read.

\subsection{PDE and PTE Permission Flags in xv6}

The x86 paging hardware provides two bits in every page-directory entry (PDE) and
page-table entry (PTE) that implement per-page isolation:

\begin{itemize}
    \item \textbf{U/S (User/Supervisor, bit 2)} --- 0 means the page is accessible only
          from ring~0 (supervisor mode); 1 means any privilege level may access the page.
          Kernel pages are mapped U/S = 0; user pages are mapped U/S = 1.
    \item \textbf{R/W (Read/Write, bit 1)} --- 0 makes the page read-only; 1 allows writes.
          Code pages are typically R/W = 0 (to catch accidental code overwrite and to
          support copy-on-write); data pages are R/W = 1.
    \item \textbf{P (Present, bit 0)} --- 1 means the page is valid; 0 causes a page fault
          on any access. The OS uses P = 0 for pages that have been swapped to disk or have
          not yet been allocated.
\end{itemize}

These bits are checked \emph{at every memory access} by the MMU hardware, with no software
involvement unless a fault occurs. The combination of separate page tables per process (for
coarse isolation) and U/S flags within the kernel's shared mappings (for fine-grained
privilege enforcement) gives xv6 a layered, robust memory isolation model.

% ============================================================
\section{CPU Isolation via Time-Slicing}
% ============================================================

\subsection{The Scheduling Abstraction}

While virtual memory isolates the \emph{memory} dimension, the CPU must also be shared
fairly and safely. The OS implements \textbf{time-slicing}: each process is given a short
\emph{time quantum} on the CPU, after which the scheduler preempts it and runs the next
runnable process. From each process's perspective, it appears to have a dedicated CPU---the
scheduling abstraction creates the \emph{illusion of parallelism}.

A \textbf{scheduler} answers three questions: \emph{who} gets the CPU, \emph{what}
(which process), and \emph{when}. The policy choice between \textbf{cooperative scheduling}
(processes voluntarily yield with a \texttt{yield()} call) and \textbf{preemptive
scheduling} (a timer interrupt forces the switch) is critical for isolation: cooperative
scheduling breaks down if any process is buggy or adversarial and never yields, whereas
preemptive scheduling guarantees that no process can monopolise the CPU indefinitely.

\subsection{CPU Isolation Concerns}

Achieving perfect CPU isolation is harder than it might appear. Even with time-slicing,
processes can interact through \emph{shared CPU state}: cache residency, TLB state,
and branch predictor state are all shared resources that time-slicing does not isolate.
These shared microarchitectural resources are precisely the attack surface exploited by
Spectre and Meltdown. For the purposes of the course, however, the scheduler model treats
the CPU as fully re-initialised between quanta.

\subsection{Time-Slicing in xv6: The Scheduler Loop}

xv6's scheduler is a simple round-robin loop that runs on each CPU:

\begin{lstlisting}[language=C, caption={proc.c --- the per-CPU scheduler loop}]
for(;;){
    sti();  // Enable interrupts so timer can fire

    acquire(&ptable.lock);
    for(p = ptable.proc; p < &ptable.proc[NPROC]; p++){
        if(p->state != RUNNABLE)
            continue;

        // Switch to chosen process.
        // It is the process's job to release ptable.lock
        // and then reacquire it before jumping back to us.
        proc = p;
        switchuvm(p);          // load process p's page table
        p->state = RUNNING;
        swtch(&cpu->scheduler, p->context);  // context switch
        switchkvm();           // restore kernel page table

        // Process is done running for now.
        proc = 0;
    }
    release(&ptable.lock);
}
\end{lstlisting}

Three isolation-relevant operations occur inside the loop. First, \texttt{switchuvm(p)}
loads CR3 with the physical address of process p's page directory, instantly replacing the
entire address space. Second, \texttt{swtch()} saves the current CPU context (register
values) and restores the saved context of process p, transferring EIP to wherever p last
left off. Third, after the process yields or is preempted and returns to the scheduler,
\texttt{switchkvm()} restores the kernel-only page table, ensuring that when the scheduler
code executes it has a clean kernel view of memory.

% ============================================================
\section{The System Call Interface}
% ============================================================

\subsection{Why a Controlled Handover?}

User code frequently needs kernel services: reading from a file, forking a child process,
allocating memory. These operations require ring~0 privileges, but user code runs at ring~3.
The transition from ring~3 to ring~0 must be \emph{controlled}---if user code could jump
to an arbitrary kernel address, it could bypass all argument validation and execute kernel
code with full privilege in an unexpected state, compromising isolation entirely.

The solution is the \textbf{trap handler}: a fixed, predetermined entry point into the
kernel that the hardware enforces. User code cannot jump to any kernel address; it can only
trigger a trap through the approved mechanism.

\subsection{The Trap Handler's Responsibilities}

When user code invokes a system call, the trap handler must perform a well-defined sequence
before executing any kernel logic:

\begin{enumerate}
    \item \textbf{Set aside the user-space context} --- save all registers onto the kernel
          stack so that the user context can be restored exactly on return.
    \item \textbf{Load a safe kernel context} --- switch to the kernel stack and kernel data
          structures. At this point the CPU is in ring~0 with the kernel's view of memory.
    \item \textbf{Sanitize arguments} --- validate every pointer and value passed by the
          user. Critically, any pointer argument must be checked to confirm it lies within
          user space. A user process must not be able to pass a kernel pointer as an argument
          to a system call and thereby trick the kernel into reading or writing kernel memory
          on its behalf.
    \item \textbf{Accept or deny the request} --- check permissions (does this process have
          access to this file?), quota limits, and any other policy requirements. Only after
          all checks pass does the kernel execute the requested operation.
\end{enumerate}

\subsection{Why Argument Sanitization Matters: The Virus Scenario}

The following code illustrates what happens when argument sanitization is absent:

\begin{lstlisting}[language=C, caption={Unsanitized system call argument enables kernel memory overwrite}]
char *buff = get_virus_code();
int fd = open("/tmp/virus", ...);
write(fd, buff, size);                         // write virus to disk
read(fd, (char *)(0x80100000), size);          // CHECK NEEDED HERE
\end{lstlisting}

The address \texttt{0x80100000} is the start of the xv6 kernel's text segment---kernel code
itself. If the \texttt{read()} system call does not verify that the destination pointer lies
within user space, the kernel will happily copy the virus payload from the file directly into
its own code pages, overwriting the instruction stream it is currently executing.
The diagram in the lecture slide shows the virus code in the user heap and also overwriting
the kernel text, connected by the blue circular arrow that represents this attack path.

A correct implementation checks that the destination address passed to \texttt{read()} falls
below the kernel boundary (\texttt{0x80000000} in xv6) before copying any data. This single
check prevents an entire class of kernel-overwrite attacks.

\subsection{Making System Calls in xv6}

On the user side, every system call is generated by a macro that follows the x86 calling
convention for software interrupts:

\begin{lstlisting}[language={[x86masm]Assembler}, caption={usys.S --- SYSCALL macro used to generate all xv6 system call stubs}]
#define SYSCALL(name)     \
  .globl name;            \
  name:                   \
      movl $SYS_##name, %eax;  \  /* put syscall number in EAX */
      int  $T_SYSCALL;          \  /* fire INT 64              */
      ret

SYSCALL(fork)
SYSCALL(exit)
\end{lstlisting}

Each system call stub places the system call number in \texttt{\%eax} (xv6's convention),
then executes \texttt{INT \$64} (\texttt{T\_SYSCALL = 64}). The \texttt{INT} instruction
atomically: saves EIP, CS, and EFLAGS onto the kernel stack; loads the CS and EIP from the
IDT gate for vector 64; and sets CPL = DPL of the gate (0 for the syscall gate). This is
the single, hardware-enforced transition from ring~3 to ring~0. The \texttt{ret} instruction
at the end of the stub is executed \emph{after} the kernel returns from handling the system
call and \texttt{IRET} restores the user context.

\subsection{Trap Numbers in xv6}

xv6 defines its trap numbers in \texttt{traps.h}. The first 20 entries are processor-defined
exceptions; xv6 adds its own starting at number 64 to avoid conflicts:

\begin{lstlisting}[language=C, caption={traps.h --- processor-defined and xv6-defined trap numbers}]
// Processor-defined:
#define T_DIVIDE   0  // divide error
#define T_DEBUG    1  // debug exception
#define T_NMI      2  // non-maskable interrupt
#define T_BRKPT    3  // breakpoint
#define T_OFLOW    4  // overflow
#define T_BOUND    5  // bounds check
#define T_ILLOP    6  // illegal opcode
#define T_DEVICE   7  // device not available
#define T_DBLFLT   8  // double fault
#define T_TSS     10  // invalid task switch segment
#define T_SEGNP   11  // segment not present
#define T_STACK   12  // stack exception
#define T_GPFLT   13  // general protection fault
#define T_PGFLT   14  // page fault
#define T_FPERR   16  // floating point error
#define T_ALIGN   17  // alignment check
#define T_MCHK    18  // machine check
#define T_SIMDERR 19  // SIMD floating point error

// xv6-defined (chosen to not overlap processor-defined):
#define T_SYSCALL 64  // system call
#define T_DEFAULT 500 // catchall
\end{lstlisting}

The choice of 64 for \texttt{T\_SYSCALL} is arbitrary but deliberate: it must be above the
processor's reserved range (0--31), and it must not conflict with any hardware interrupt
that has been mapped by the PIC (Programmable Interrupt Controller) into the IDT. The
\texttt{T\_DEFAULT} catchall at 500 exists so that unexpected interrupts do not silently
execute garbage code.

\subsection{Scale of System Call Interfaces}

The system call table is the entire visible surface area between user code and the kernel.
The size of this table is a design choice that reflects the OS's philosophy:

\begin{center}
\begin{tabular}{ll}
\toprule
\textbf{OS} & \textbf{System Calls} \\
\midrule
xv6         & 21 \\
Linux       & 300+ \\
FreeBSD     & 500+ \\
Windows 10--11 & $\approx$465 core / 2000+ total \\
\bottomrule
\end{tabular}
\end{center}

xv6's minimalist table of 21 syscalls is intentional: the course uses xv6 precisely because
it implements all core OS mechanisms without the complexity of a production system. A useful
debugging tool on Linux is \texttt{strace}, which traces every system call made by a
process---it is invaluable for understanding what a program is actually doing at the
OS boundary.

% ============================================================
\section{Isolation: A Solved Problem?}
\label{sec:limits}
% ============================================================

A natural question after seeing all the isolation mechanisms above is: is isolation a
solved problem? The answer, unfortunately, is no. Two landmark security vulnerabilities
disclosed in January 2018 demonstrated that hardware-level attacks can break the isolation
guarantees described in this lecture, even when the software is perfectly correct.

\textbf{Meltdown} exploits the fact that modern processors speculatively execute
instructions before checking whether those instructions are permitted. A user process can
cause the CPU to speculatively load a kernel memory address (which would ordinarily be
blocked by the U/S page-table bit), and before the fault is raised, use the speculatively
loaded value to encode information into the CPU cache via a carefully crafted cache-timing
side channel. By subsequently measuring cache access times, the attacker recovers the kernel
memory contents without ever having architectural access to them. Meltdown was described by
its discoverers as something that ``breaks the most fundamental isolation between user
applications and the operating system.''

\textbf{Spectre} exploits branch prediction and speculative execution to cause a victim
process to speculatively access its own memory in a way that leaves a cache-timing side
channel readable by an attacker in a different process. Unlike Meltdown, Spectre does not
require accessing kernel memory; it works across process boundaries. Its discoverers
characterised it as breaking ``the isolation between different applications.''

Both attacks are deeply rooted in microarchitectural optimisations that were designed for
performance, not isolation. They underscore the lecture's key point about the safety/usability
tradeoff and the ``performance'' axis being a question mark: the very optimisations that make
modern CPUs fast (speculative execution, out-of-order execution, branch prediction) create
new attack surfaces that can erode even well-designed software isolation. Mitigations
(KPTI, retpoline, microcode updates) have been deployed, but they come at measurable
performance cost and do not fully eliminate the underlying problem.

% ============================================================
\section{Summary}
% ============================================================

\begin{itemize}
    \item \textbf{Isolation definition}: the property by which an OS entity is protected
          from the faults, misbehaviour, or malice of other entities sharing the same hardware.
    \item \textbf{Benefits}: failure isolation (faults don't propagate), performance
          isolation (no starvation), correctness isolation (results unaffected by other programs).
    \item \textbf{Costs}: 100\% isolation is useless---a process must communicate to be
          useful. We want \emph{some} isolation \emph{sometimes}, calibrated to workload and
          trust relationships.
    \item \textbf{Necessary conditions}: isolation is appropriate when resources are shared
          \emph{and} entities are non-cooperative.
    \item \textbf{Tradeoffs}: stronger isolation $\to$ safer but less usable; the goal of
          OS design is to push the Pareto frontier toward both high safety and high usability.
    \item \textbf{Mechanisms}: hardware protection rings, virtual address spaces, time-slicing,
          system call interface, and user/group permissions.
    \item \textbf{Protection rings}: x86 has four rings; practical OSes use only rings 0 and 3.
          CPL (lower 2 bits of \texttt{\%cs}) determines current privilege. Ring~0 guards
          CR3, CS writes, memory privilege checks, I/O ports, and control registers.
    \item \textbf{Segment selectors}: encode GDT index, TI (GDT/LDT), and RPL. CPL is
          derived from the RPL of the loaded \texttt{\%cs} selector.
    \item \textbf{Memory isolation}: each process gets its own virtual address space via
          a private page table. U/S and R/W bits in PDE/PTE entries enforce per-page
          access control at hardware speed.
    \item \textbf{xv6 GDT}: kernel segments have DPL = 0; user segments have DPL = 3
          (\texttt{DPL\_USER}). All four segments are flat (base 0, full 4~GB limit);
          isolation comes from DPL and from paging, not from segment limits.
    \item \textbf{CPU isolation}: time-slicing via a preemptive scheduler; xv6 implements
          round-robin using \texttt{switchuvm()} + \texttt{swtch()} + \texttt{switchkvm()}.
    \item \textbf{System call interface}: the only approved ring-3 $\to$ ring-0 pathway;
          user code fires \texttt{INT \$64}; the trap handler saves context, loads kernel
          context, \emph{sanitizes arguments}, and accepts or denies the request.
    \item \textbf{Argument sanitization}: every pointer passed to a syscall must be
          verified to lie within user space; failure to do so can allow a process to overwrite
          kernel memory.
    \item \textbf{Limits of isolation}: Meltdown and Spectre (2018) showed that hardware
          speculative execution can break software isolation at a microarchitectural level,
          even when all software checks are correct.
\end{itemize}

% ============================================================
\section{References}
% ============================================================

\begin{itemize}
    \item Tumanov, A. \textit{CS3210: Design of Operating Systems --- Lecture 4 Slides}.
          Georgia Institute of Technology. January 27, 2026.
    \item Cox, R.; Kaashoek, F.; Morris, R. \textit{xv6: a simple, Unix-like teaching
          operating system}. MIT CSAIL. \url{https://pdos.csail.mit.edu/6.828/2022/xv6.html}
    \item Lipp, M., et al. ``Meltdown: Reading Kernel Memory from User Space.''
          \textit{USENIX Security Symposium}. 2018. \url{https://meltdownattack.com/}
    \item Kocher, P., et al. ``Spectre Attacks: Exploiting Speculative Execution.''
          \textit{IEEE Security \& Privacy}. 2019.
          \url{https://googleprojectzero.blogspot.com/2018/01/reading-privileged-memory-with-side.html}
    \item Boos, K., et al. ``Theseus: an Experiment in Operating System Structure and State
          Management.'' \textit{USENIX OSDI}. 2020.
          \url{https://www.usenix.org/conference/osdi20/presentation/boos}
    \item Intel Corporation. \textit{Intel 64 and IA-32 Architectures Software Developer's
          Manual, Volume 3A: System Programming Guide}.
          \url{https://www.intel.com/content/www/us/en/developer/articles/technical/intel-sdm.html}
    \item Silberschatz, A.; Galvin, P.~B.; Gagne, G. \textit{Operating System Concepts},
          10th ed. John Wiley \& Sons. 2018.
\end{itemize}

% ============================================================
\section*{Personal Notes}
% ============================================================
% Add your own notes below this line

\end{document}
